\chapter[FORMAT DES DONNÉES MIC-E]{Format des données Mic-E}\label{format-des-donnees-mic-e}

\begingroup
\setlength{\arrayrulewidth}{0.6pt}
\renewcommand{\lit}[1]{\begingroup\setlength{\fboxsep}{1pt}\colorbox{APRSMarkerYellow}{\strut\bfseries\aprsorigtexttt{#1}}\endgroup}

\newcommand{\chTenToc}[1]{\phantomsection\addcontentsline{toc}{section}{#1}}
\newcommand{\chTenHead}[1]{{\sffamily\bfseries\itshape #1}}
\newcommand{\chTenGray}{\cellcolor{gray!20}}
\newcommand{\chTenBlack}[1]{\cellcolor{black}\textcolor{white}{\bfseries #1}}
\newfontfamily{\chTenSpaceFont}{FreeSans}
\newcommand{\chTenSpace}{\begingroup\chTenSpaceFont ⎵\endgroup}
\newcommand{\chTenLitSpace}{\begingroup\setlength{\fboxsep}{1pt}\colorbox{APRSMarkerYellow}{\strut\bfseries\chTenSpaceFont ⎵}\endgroup}
\newcommand{\chTenItems}{\setlength{\itemsep}{0.35em}\setlength{\parskip}{0pt}}

\chTenToc{Format des données Mic-E}
\begin{aprssideblock}{Format des données\\Mic-E}
Dans le format de données Mic-E, la position, le cap, la vitesse et le symbole
d'affichage de la station, ainsi qu'un chemin de digipeaters APRS et un
commentaire de position Mic-E, sont regroupés dans les champs Destination
Address et Information AX.25.

Le champ Information peut également contenir, en option, un statut Mic-E. Le
statut Mic-E peut contenir le locator Maidenhead et l'altitude de la station.

Les paquets Mic-E peuvent être très courts. Au minimum, sans indicatifs dans le
champ Digipeater Addresses et sans données de télémétrie ni texte de statut
Mic-E optionnels, un paquet Mic-E complet ne comporte que 25 octets (hors FCS
et flags).

Le format de données Mic-E n'est pas seulement utilisé dans le
\href{https://web.tapr.org/product_docs/Mic-E/mic-edev/mic-e.manual.pdf}{Microphone Encoder} ;
il l'est aussi dans le
\href{https://www.tapr.org/pdf/DCC1999-APRS-MicLite-WB4APR.pdf}{PIC Encoder} et
dans les radios Kenwood et Yaesu.
\end{aprssideblock}

\chTenToc{Charge utile Mic-E}
\begin{aprssideblock}{Charge utile Mic-E}
Le format de données Mic-E permet de transporter une grande quantité de
données dans un paquet très court. Les données sont réparties entre le champ
Destination Address et le champ Information d'une UI-frame AX.25 standard.

\textbf{Champ Destination Address} --- Le champ Destination Address de
7 octets contient les informations encodées suivantes :

\begin{itemize}\chTenItems
\item les 6 chiffres de latitude ;
\item un commentaire de position Mic-E sur 3 bits, spécifiant l'un des
  7 commentaires de position Standard, l'un des 7 commentaires de position
  Custom ou un commentaire de position Emergency ;
\item les indicateurs North/South et West/East ;
\item l'indicateur Longitude Offset ;
\item le code de chemin de digipeaters APRS générique.
\end{itemize}

Bien que la Destination Address paraisse assez inhabituelle, elle reste une
adresse AX.25 valide, constituée uniquement de valeurs ASCII 7 bits imprimables
(décalées d'un bit vers la gauche) --- voir la spécification
\textit{Amateur Packet-Radio Link-Layer Protocol} pour tous les détails sur le
format des adresses AX.25 standard.

\textbf{Champ Information} --- Ce champ contient les données suivantes :

\begin{itemize}\chTenItems
\item la longitude encodée ;
\item le cap et la vitesse encodés ;
\item le Symbol Code d'affichage et le Symbol Table Identifier ;
\item un champ optionnel contenant soit des données de télémétrie Mic-E
  (désormais obsolètes), soit une chaîne de texte de statut Mic-E. La chaîne de
  statut peut contenir du texte en clair, un locator Maidenhead ou l'altitude
  de la station.
\end{itemize}
\end{aprssideblock}

\chTenToc{Champ Destination Address Mic-E}
\begin{aprssideblock}{Champ Destination\\Address Mic-E}
Le champ Destination Address AX.25 standard comporte 7 octets, contenant
6 caractères d'indicatif et le SSID (ainsi qu'un certain nombre d'autres bits
qui ne nous intéressent pas ici). Lorsqu'il est utilisé pour transporter des
données Mic-E, ce champ présente toutefois un format très différent :
\end{aprssideblock}

\begin{center}
\setlength{\tabcolsep}{2pt}
\renewcommand{\arraystretch}{1.75}
\footnotesize\sffamily
\hspace*{15pt}%
\begin{tabular}[t]{@{}|*{7}{>{\centering\arraybackslash}p{0.1215\textwidth}|}@{}}
\hline
\multicolumn{7}{|l|}{\cellcolor{black}\textcolor{white}{\sffamily\bfseries\itshape
Mic-E Data --- DESTINATION ADDRESS FIELD Format}} \\
\hline
\chTenHead{Lat Digit 1\newline + Bit A} &
\chTenHead{Lat Digit 2\newline + Bit B} &
\chTenHead{Lat Digit 3\newline + Bit C} &
\chTenHead{Lat Digit 4\newline + N/S Lat\newline Indicator} &
\chTenHead{Lat Digit 5\newline + Longitude\newline Offset} &
\chTenHead{Lat Digit 6\newline + W/E Long\newline Indicator} &
\chTenHead{APRS\newline Digi Path\newline Code} \\
\hline
\multicolumn{1}{|c|}{\makebox[0pt][r]{Bytes:\hspace{37pt}}1} & 1 & 1 & 1 & 1 & 1 & 1 \\
\hline
\end{tabular}
\end{center}

\begin{aprsbodyblock}
Le champ Destination Address contient :

\begin{itemize}\chTenItems
\item six chiffres de latitude encodés spécifiant les degrés (chiffres 1 et 2),
  les minutes (chiffres 3 et 4) et les centièmes de minute (chiffres 5 et 6) ;
\item un commentaire de position Mic-E sur 3 bits (bits A, B et C) ;
\item l'indicateur de latitude North/South ;
\item le Longitude Offset (qui ajoute 0 degré ou 100 degrés au calcul de la
  longitude dans le champ Information) ;
\item l'indicateur de longitude West/East ;
\item le chemin de digipeaters APRS générique (encodé dans le SSID).
\end{itemize}
\end{aprsbodyblock}

\chTenToc{Encodage du champ Destination Address}
\begin{aprssideblock}{Encodage du champ\\Destination Address}
Le tableau de la page suivante présente l'encodage des 6 premiers octets du
champ Destination Address, pour toutes les combinaisons du chiffre de latitude,
du commentaire de position Mic-E sur 3 bits (A/B/C), des indicateurs de
latitude/longitude et du Longitude Offset.

L'encodage prend en charge l'ambiguïté de position.

Les caractères ASCII indiqués dans le tableau sont décalés d'un bit vers la
gauche avant transmission.
\end{aprssideblock}

\needspace{22\baselineskip}
\begin{center}
\sffamily\bfseries Encodage du champ Destination Address Mic-E (octets 1--6)

\medskip
\footnotesize\sffamily\normalfont
\setlength{\tabcolsep}{2pt}
\renewcommand{\arraystretch}{1.33}
\newcommand{\desthead}{%
\rule{0pt}{30pt}\chTenHead{\shortstack{ASCII\\Char}} &
\chTenHead{\shortstack{Lat\\Digit}} &
\chTenHead{\shortstack{Position \\Comment\\A/B/C}} & \chTenHead{N/S} &
\chTenHead{\shortstack{Long\\Offset}} & \chTenHead{W/E}}
\makebox[\textwidth][c]{\hspace*{9pt}%
\begin{tabular}[t]{@{}|>{\columncolor{gray!20}\bfseries\centering\arraybackslash}p{31.5pt}|>{\centering\arraybackslash}p{31.5pt}|>{\centering\arraybackslash}p{51.5pt}|>{\centering\arraybackslash}p{31.5pt}|>{\centering\arraybackslash}p{38pt}|>{\centering\arraybackslash}p{32pt}|@{}}
\cline{2-6}
\multicolumn{1}{r|}{Byte:} & 1--6 & 1--3 & 4 & 5 & 6 \\
\hline\rowcolor{gray!20}\desthead\\\hline
0&0&0&South&+0&East\\\hline
1&1&0&South&+0&East\\\hline
2&2&0&South&+0&East\\\hline
3&3&0&South&+0&East\\\hline
4&4&0&South&+0&East\\\hline
5&5&0&South&+0&East\\\hline
6&6&0&South&+0&East\\\hline
7&7&0&South&+0&East\\\hline
8&8&0&South&+0&East\\\hline
9&9&0&South&+0&East\\\hline
A&0&1 (Custom)&\chTenGray&\chTenGray&\chTenGray\\\hline
B&1&1 (Custom)&\chTenGray&\chTenGray&\chTenGray\\\hline
C&2&1 (Custom)&\chTenGray&\chTenGray&\chTenGray\\\hline
D&3&1 (Custom)&\chTenGray&\chTenGray&\chTenGray\\\hline
E&4&1 (Custom)&\chTenGray&\chTenGray&\chTenGray\\\hline
F&5&1 (Custom)&\chTenGray&\chTenGray&\chTenGray\\\hline
G&6&1 (Custom)&\chTenGray&\chTenGray&\chTenGray\\\hline
\end{tabular}
\hspace{10.5pt}%
\begin{tabular}[t]{@{}|>{\columncolor{gray!20}\bfseries\centering\arraybackslash}p{31.5pt}|>{\centering\arraybackslash}p{31.5pt}|>{\centering\arraybackslash}p{51.5pt}|>{\centering\arraybackslash}p{31.5pt}|>{\centering\arraybackslash}p{38pt}|>{\centering\arraybackslash}p{32pt}|@{}}
\cline{2-6}
\multicolumn{1}{r|}{Byte:} & 1--6 & 1--3 & 4 & 5 & 6 \\
\hline\rowcolor{gray!20}\desthead\\\hline
H&7&1 (Custom)&\chTenGray&\chTenGray&\chTenGray\\\hline
I&8&1 (Custom)&\chTenGray&\chTenGray&\chTenGray\\\hline
J&9&1 (Custom)&\chTenGray&\chTenGray&\chTenGray\\\hline
K&space&1 (Custom)&\chTenGray&\chTenGray&\chTenGray\\\hline
L&space&0&South&+0&East\\\hline
P&0&1 (Std)&North&+100&West\\\hline
Q&1&1 (Std)&North&+100&West\\\hline
R&2&1 (Std)&North&+100&West\\\hline
S&3&1 (Std)&North&+100&West\\\hline
T&4&1 (Std)&North&+100&West\\\hline
U&5&1 (Std)&North&+100&West\\\hline
V&6&1 (Std)&North&+100&West\\\hline
W&7&1 (Std)&North&+100&West\\\hline
X&8&1 (Std)&North&+100&West\\\hline
Y&9&1 (Std)&North&+100&West\\\hline
Z&space&1 (Std)&North&+100&West\\\hline
\end{tabular}
\hspace*{-9pt}}
\end{center}

\begin{aprsbodyblock}
\textbf{Note :} les caractères ASCII \lit{A}--\lit{K} ne sont pas utilisés
dans les octets d'adresse 4 à 6.

Par exemple, pour une station située à une latitude de 33 degrés
25,64 minutes nord, dans l'hémisphère ouest, avec un Longitude Offset de
+0 degré, et transmettant les bits de commentaire de position standard 1/0/0,
l'encodage des 6 premiers octets du champ Destination Address est le suivant :
\end{aprsbodyblock}

\begin{center}
\small\sffamily
\setlength{\tabcolsep}{3pt}\renewcommand{\arraystretch}{1.75}
\setlength{\arrayrulewidth}{0.75pt}
\begin{tabular}[t]{@{}|
  >{\columncolor{gray!20}\raggedleft\arraybackslash}p{65pt}|
  *{6}{>{\centering\arraybackslash}p{43.5pt}|}
  @{}}
\hhline{~|------|}
\multicolumn{1}{r|}{\makebox[0pt][r]{Destination Address Byte:\hspace{6pt}}} & \chTenGray\bfseries\itshape 1&\chTenGray\bfseries\itshape 2&\chTenGray\bfseries\itshape 3&\chTenGray\bfseries\itshape 4&\chTenGray\bfseries\itshape 5&\chTenGray\bfseries\itshape 6\\\hline
\chTenHead{Latitude Digit}&3&3&2&5&6&4\\\hline
\chTenHead{Position Comment}&\vspace{3pt}\shortstack{Bit A\\= 1 (Std)}&\vspace{3pt}\shortstack{Bit B\\= 0}&\vspace{3pt}\shortstack{Bit C\\= 0}&\chTenGray&\chTenGray&\chTenGray\\\hline
\chTenHead{N/S Indicator}&\chTenGray&\chTenGray&\chTenGray&North&\chTenGray&\chTenGray\\\hline
\chTenHead{Long Offset}&\chTenGray&\chTenGray&\chTenGray&\chTenGray&+0&\chTenGray\\\hline
\chTenHead{W/E Indicator}&\chTenGray&\chTenGray&\chTenGray&\chTenGray&\chTenGray&West\\\hline
\chTenHead{Dest Address\newline (ASCII Char)}&\ttfamily S&\ttfamily 3&\ttfamily 2&\ttfamily U&\ttfamily 6&\ttfamily T\\\hline
\end{tabular}
\end{center}

\chTenToc{Commentaires de position Mic-E}
\begin{aprssideblock}{Commentaires de position Mic-E}
Les trois premiers octets du champ Destination Address contiennent trois bits
de commentaire de position : A, B et C. Ces bits permettent de spécifier l'un
des 15 types de commentaire de position :

\begin{itemize}\chTenItems
\item 7 commentaires de position Standard ;
\item 7 commentaires de position Custom ;
\item 1 commentaire de position Emergency.
\end{itemize}

Pour les 7 commentaires de position Standard, un ou plusieurs bits de
commentaire de position valent \lit{1}, indiqués par \textit{1 (Std)} dans le
tableau d'encodage du champ Destination Address Mic-E.

Pour les 7 commentaires de position Custom, un ou plusieurs bits de
commentaire de position valent \lit{1}, indiqués par \textit{1 (Custom)} dans ce
même tableau.

Pour le commentaire de position Emergency, les trois bits de commentaire de
position valent \lit{0}.

Le tableau suivant présente l'encodage des types de commentaire de position
Mic-E pour toutes les combinaisons des bits de commentaire de position A/B/C :

\par\medskip
\sffamily\bfseries Types de commentaires de position Mic-E

\smallskip
\sffamily\normalfont
\setlength{\tabcolsep}{2.5pt}\renewcommand{\arraystretch}{1.35}
\begin{tabular}[t]{!{\vrule width 1pt}>{\columncolor{gray!20}\centering\arraybackslash}p{15pt}|>{\columncolor{gray!20}\centering\arraybackslash}p{15pt}|>{\columncolor{gray!20}\centering\arraybackslash}p{15pt}|>{\raggedright\arraybackslash}p{0.23\linewidth}|>{\raggedright\arraybackslash}p{0.23\linewidth}|}\hline
\rowcolor{gray!20}\chTenHead{A}&\chTenHead{B}&\chTenHead{C}&\chTenHead{Commentaire de position\newline Standard}&\chTenHead{Commentaire de position\newline Custom}\\\hline
1&1&1&M0: Off Duty&C0: Custom-0\\\hline
1&1&0&M1: En Route&C1: Custom-1\\\hline
1&0&1&M2: In Service&C2: Custom-2\\\hline
1&0&0&M3: Returning&C3: Custom-3\\\hline
0&1&1&M4: Committed&C4: Custom-4\\\hline
0&1&0&M5: Special&C5: Custom-5\\\hline
0&0&1&M6: Priority&C6: Custom-6\\\hline
0&0&0&\multicolumn{2}{c|}{Emergency}\\\hline
\end{tabular}

Les commentaires de position Standard et le commentaire de position Emergency
ont la même signification pour toutes les stations APRS. Les commentaires de
position Custom peuvent recevoir n'importe quelle signification arbitraire.

\textbf{Note :} la prise en charge des commentaires de position Custom est
facultative. Les unités Mic-E d'origine ne prennent pas en charge les
commentaires de position Custom.

\textbf{Note :} si les bits de commentaire de position A/B/C contiennent un
mélange de \lit{1} Standard et de \lit{1} Custom, le type de commentaire de
position est « unknown ».
\end{aprssideblock}

\begin{aprsbodyblock}
Quelques exemples d'encodage du type de commentaire de position :

\begin{center}
\small\sffamily
\setlength{\tabcolsep}{3.5pt}\renewcommand{\arraystretch}{1.70}
\makebox[\linewidth][c]{\hspace*{-13.5pt}%
\begin{tabular}[t]{|>{\columncolor{gray!20}\centering\arraybackslash}p{0.22\linewidth}|>{\centering\arraybackslash}p{0.31\linewidth}|>{\centering\arraybackslash}p{0.14\linewidth}|>{\raggedright\arraybackslash}p{0.18\linewidth}|}\hline
\rowcolor{gray!20}
\chTenHead{3 premiers octets de la\newline Destination Address}&
\chTenHead{Bits de commentaire de\newline position A/B/C}&
\chTenHead{Type de\newline commentaire de position}&\chTenHead{Signification}\\\hline
\ttfamily S32&Standard 1 / 0 / 0&Standard&M3: Returning\\\hline
\ttfamily F2D&Custom 1 / 0 / Custom 1&Custom&C2: Custom-2\\\hline
\ttfamily 234&0 / 0 / 0&Emergency&Emergency\\\hline
\end{tabular}%
\hspace*{13.5pt}}
\end{center}
\end{aprsbodyblock}

\chTenToc{Champ SSID de Destination Address}
\begin{aprssideblock}{Champ SSID de\\Destination Address}
Le SSID du champ Destination Address d'un paquet Mic-E est codé pour spécifier
soit un chemin de digipeaters VIA conventionnel (contenu dans le champ
Digipeater Addresses de la trame AX.25), soit l'un des 15 chemins de
digipeaters APRS génériques. Voir le chapitre 4 : Données APRS dans les champs
Destination Address et Source Address AX.25.
\end{aprssideblock}

\chTenToc{Champ Information Mic-E}
\begin{aprssideblock}{Champ Information\\Mic-E}
Le champ Information sert à compléter le Position Report commencé dans le
champ Destination Address. L'encodage utilisé diffère de celui de la
Destination Address, car le contenu n'est pas contraint à de l'ASCII 7 bits
imprimable décalé, comme il l'est dans l'adresse. Toutefois, le binaire 8 bits
complet n'est pas utilisé : toutes les valeurs sont décalées de 28 et d'autres
opérations (décrites ci-dessous) sont effectuées sur certaines valeurs afin de
rendre presque toutes les données imprimables en ASCII.

Le format du champ Information est le suivant :
\end{aprssideblock}

\begin{center}
\footnotesize\sffamily\setlength{\tabcolsep}{1.5pt}\renewcommand{\arraystretch}{1.40}
\makebox[\textwidth][c]{\hspace*{11.5pt}%
\begin{tabular}[t]{@{}|>{\centering\arraybackslash}p{31.5pt}|*{3}{>{\centering\arraybackslash}p{32pt}|}*{3}{>{\centering\arraybackslash}p{34.5pt}|}>{\centering\arraybackslash}p{39pt}|>{\centering\arraybackslash}p{32pt}|>{\centering\arraybackslash}p{110.5pt}|@{}}\hline
\multicolumn{10}{|l|}{\cellcolor{gray!20}\sffamily\bfseries\itshape Mic-E Data --- INFORMATION FIELD Format}\\\hline
\multirow{2}{31.5pt}{\centering\chTenHead{Data\newline Type\newline ID}}&
\multicolumn{3}{c|}{\chTenHead{Longitude}}&
\multicolumn{3}{c|}{\chTenHead{Speed and Course}}&
\multirow{2}{39pt}{\centering\chTenHead{Symbol\newline Code}}&
\multirow{2}{32pt}{\centering\chTenHead{Sym\newline Table\newline ID}}&
\cellcolor{gray!20}\chTenHead{Mic-E Telemetry Data (obsolete)}\\
\hhline{|~|---|---|~|~|-|}
&\ttfamily d+28&\ttfamily m+28&\ttfamily h+28&\ttfamily SP+28&\ttfamily DC+28&\ttfamily SE+28&&&\cellcolor{gray!20}\chTenHead{Mic-E Status Text}\\\hline
\makebox[0pt][r]{Bytes:\hspace{1.2em}}\makebox[\linewidth][c]{1}&1&1&1&1&1&1&1&1&n\\\hline
\end{tabular}%
\hspace*{-11.5pt}}
\end{center}

\chTenToc{Données du champ Information}
\begin{aprssideblock}{Données du champ\\Information}
Les 9 premiers octets du champ Information contiennent l'APRS Data Type
Identifier, la longitude, la vitesse, le cap et les données de symbole.

L'APRS Data Type Identifier est l'un des suivants :

\begin{tabular}[t]{@{}>{\ttfamily}p{0.07\linewidth}p{0.86\linewidth}@{}}
\lit{`} & données GNSS actuelles (mais non utilisées dans les radios Kenwood
TM-D700) ; \\
\lit{'} & anciennes données GNSS (ou données GNSS actuelles dans les radios
Kenwood TM-D700) ; \\
\texttt{0x1c} & données GNSS actuelles (obsolète --- unités beta Rev. 0
uniquement) ; \\
\texttt{0x1d} & anciennes données GNSS (obsolète --- unités beta Rev. 0
uniquement).
\end{tabular}

\textbf{\underline{NOTE IMPORTANTE :}} comme expliqué en détail ci-dessous, certains
octets du champ Information sont des caractères ASCII non imprimables.
L'utilisation de caractères non imprimables rend le dépannage plus difficile.
Dans certaines circonstances (configuration incorrecte du TNC ou logiciel
inapproprié, par exemple), certains de ces caractères non imprimables peuvent
être supprimés, ce qui fait paraître le champ Information plus court qu'il ne
l'est réellement. Cela entraîne un décodage incorrect. En particulier, si le
champ Information paraît comporter moins de 9 octets, le paquet doit être
ignoré.
\end{aprssideblock}

\chTenToc{Encodage des degrés de longitude}
\begin{aprssideblock}{Encodage des degrés\\de longitude}
L'octet \lit{d+28} du champ Information contient la valeur encodée des degrés
de longitude, dans la plage 0--179 degrés.

(Noter que, pour les longitudes comprises entre 0 et 9 degrés, le Longitude
Offset est de +100 degrés) :

\par\smallskip
\sffamily\bfseries Encodage Mic-E des degrés de longitude

\smallskip
\footnotesize\sffamily\normalfont
\setlength{\tabcolsep}{5pt}\renewcommand{\arraystretch}{1.61}
\hspace*{-7pt}%
\begin{tabular}[t]{|>{\columncolor{gray!20}\bfseries\centering\arraybackslash}p{32pt}|>{\centering\arraybackslash}p{25pt}|>{\centering\arraybackslash}p{25.5pt}|>{\centering\arraybackslash}p{32pt}|}\hline
\rowcolor{gray!20}\rule{0pt}{25pt}\chTenHead{\shortstack{Long\\Deg}}&\chTenHead{\shortstack{ASCII\\Char}}&\chTenHead{d+28}&\chTenHead{\shortstack{Long\\Offset}}\\\hline
0&v&118&\lit{+100}\\\hline
1&w&119&\lit{+100}\\\hline
2&x&120&\lit{+100}\\\hline
3&y&121&\lit{+100}\\\hline
4&z&122&\lit{+100}\\\hline
5&\{&123&\lit{+100}\\\hline
6&|&124&\lit{+100}\\\hline
7&\}&125&\lit{+100}\\\hline
8&\textasciitilde&126&\lit{+100}\\\hline
9&\chTenBlack{DEL}&127&\lit{+100}\\\hline
10&\& &38&+0\\\hline
11&\aprsorigtexttt{'}&39&+0\\\hline
12&(&40&+0\\\hline
\ldots&\ldots&\ldots&\ldots\\\hline
97&\}&125&+0\\\hline
98&\textasciitilde&126&+0\\\hline
99&\chTenBlack{DEL}&127&+0\\\hline
\end{tabular}\hspace{35pt}%
\begin{tabular}[t]{|>{\columncolor{gray!20}\bfseries\centering\arraybackslash}p{32pt}|>{\centering\arraybackslash}p{25pt}|>{\centering\arraybackslash}p{25.5pt}|>{\centering\arraybackslash}p{32pt}|}\hline
\rowcolor{gray!20}\rule{0pt}{25pt}\chTenHead{\shortstack{Long\\Deg}}&\chTenHead{\shortstack{ASCII\\Char}}&\chTenHead{d+28}&\chTenHead{\shortstack{Long\\Offset}}\\\hline
100&l&108&+100\\\hline
101&m&109&+100\\\hline
102&n&110&+100\\\hline
103&o&111&+100\\\hline
104&p&112&+100\\\hline
105&q&113&+100\\\hline
106&r&114&+100\\\hline
107&s&115&+100\\\hline
108&t&116&+100\\\hline
109&u&117&+100\\\hline
110&\& &38&+100\\\hline
111&\aprsorigtexttt{'}&39&+100\\\hline
112&(&40&+100\\\hline
\ldots&\ldots&\ldots&\ldots\\\hline
177&i&105&+100\\\hline
178&j&106&+100\\\hline
179&k&107&+100\\\hline
\end{tabular}

\normalsize
Le tableau montre que l'encodage est divisé en quatre parties distinctes :

\begin{itemize}\chTenItems
\item \underline{0--9 degrés} : \lit{d+28} est compris entre 118 et 127 en
  décimal, ce qui correspond aux caractères ASCII \lit{v} à \lit{DEL}.

  \textbf{Note importante :} le Longitude Offset est réglé sur +100 degrés
  lorsque la longitude est comprise entre 0 et 9 degrés.
\item \underline{10--99 degrés} : \lit{d+28} est compris entre 38 et 127 en
  décimal (ce qui correspond aux caractères ASCII \lit{\&} à \lit{DEL}) et le Longitude
  Offset est de +0 degré.
\item \underline{100--109 degrés} : \lit{d+28} est compris entre 108 et 117 en
  décimal, ce qui correspond aux caractères ASCII \lit{l} (lettre « L »
  minuscule) à \lit{u}, et le Longitude Offset est de +100 degrés.
\item \underline{110--179 degrés} : \lit{d+28} est compris entre 38 et 107 en
  décimal (ce qui correspond aux caractères ASCII \lit{\&} à \lit{k}) et le
  Longitude Offset est de +100 degrés.
\end{itemize}

La plage globale des valeurs valides de \lit{d+28} est donc 38--127 en décimal
(ce qui correspond aux caractères ASCII \lit{\&} à \lit{DEL}).

Tous ces caractères, \underline{à l'exception de \lit{DEL} pour 9 et 99
degrés}, sont des caractères ASCII imprimables.

Pour décoder la valeur des degrés de longitude :

\begin{enumerate}\chTenItems
\item soustraire 28 de la valeur \lit{d+28} pour obtenir \lit{d} ;
\item si le Longitude Offset est de +100 degrés, ajouter 100 à \lit{d} ;
\item soustraire 80 si $180 \le \lit{d} \le 189$ (c'est-à-dire si la longitude est
  comprise entre 100 et 109 degrés) ;
\item sinon, soustraire 190 si $190 \le \lit{d} \le 199$ (c'est-à-dire si la
  longitude est comprise entre 0 et 9 degrés).
\end{enumerate}
\end{aprssideblock}

\chTenToc{Encodage des minutes de longitude}
\begin{aprssideblock}{Encodage des minutes\\de longitude}
L'octet \lit{m+28} du champ Information contient la valeur encodée des minutes
de longitude, dans la plage 0--59 minutes :

\par\smallskip
\sffamily\bfseries Encodage Mic-E des minutes de longitude

\smallskip
\footnotesize\sffamily\normalfont
\setlength{\tabcolsep}{4.5pt}\renewcommand{\arraystretch}{1.60}
\hspace*{-3pt}%
\begin{tabular}[t]{|>{\columncolor{gray!20}\bfseries\centering\arraybackslash}p{33pt}|>{\centering\arraybackslash}p{26pt}|>{\centering\arraybackslash}p{26.5pt}|}\hline
\rowcolor{gray!20}\rule{0pt}{25pt}\chTenHead{\shortstack{Long\\Mins}}&\chTenHead{\shortstack{ASCII\\Char}}&\chTenHead{m+28}\\\hline
0&X&88\\\hline
1&Y&89\\\hline
2&Z&90\\\hline
3&[&91\\\hline
4&\textbackslash&92\\\hline
5&]&93\\\hline
6&\textasciicircum&94\\\hline
7&\_&95\\\hline
8&\aprsorigtexttt{`}&96\\\hline
9&a&97\\\hline
\end{tabular}\hspace{28pt}%
\begin{tabular}[t]{|>{\columncolor{gray!20}\bfseries\centering\arraybackslash}p{33pt}|>{\centering\arraybackslash}p{26pt}|>{\centering\arraybackslash}p{26.5pt}|}\hline
\rowcolor{gray!20}\rule{0pt}{25pt}\chTenHead{\shortstack{Long\\Mins}}&\chTenHead{\shortstack{ASCII\\Char}}&\chTenHead{m+28}\\\hline
10&\& &38\\\hline
11&\aprsorigtexttt{'}&39\\\hline
12&(&40\\\hline
13&)&41\\\hline
14&*&42\\\hline
\ldots&&\\\hline
56&T&84\\\hline
57&U&85\\\hline
58&V&86\\\hline
59&W&87\\\hline
\end{tabular}

\normalsize
Le tableau montre que l'encodage est divisé en deux parties distinctes :

\begin{itemize}\chTenItems
\item \underline{0--9 minutes} : \lit{m+28} est compris entre 88 et 97 en
  décimal, ce qui correspond aux caractères ASCII \lit{X} à \lit{a} ;
\item \underline{10--59 minutes} : \lit{m+28} est compris entre 38 et 87 en
  décimal (ce qui correspond aux caractères ASCII \lit{\&} à \lit{W}).
\end{itemize}

La plage globale des valeurs valides de \lit{m+28} est donc 38--97 en décimal
(ce qui correspond aux caractères ASCII \lit{\&} à \lit{a}). Tous ces
caractères sont des caractères ASCII imprimables.

Pour décoder la valeur des minutes de longitude :

\begin{enumerate}\chTenItems
\item soustraire 28 de la valeur \lit{m+28} pour obtenir \lit{m} ;
\item soustraire 60 si $\lit{m} \ge 60$ (c'est-à-dire si les minutes de longitude
  sont comprises entre 0 et 9).
\end{enumerate}
\end{aprssideblock}

\chTenToc{Encodage des centièmes de minute de longitude}
\begin{aprssideblock}{Encodage des\\centièmes de minute\\de longitude}
L'octet \lit{h+28} du champ Information contient la valeur encodée des
centièmes de minute de longitude, dans la plage 0--99 minutes. Cet octet
prend une valeur comprise entre 28 en décimal (correspondant à 0 centième de
minute) et 127 en décimal (correspondant à 99 centièmes de minute).

Pour décoder la valeur des centièmes de minute de longitude, soustraire 28 de
la valeur \lit{h+28}.

Toutes les valeurs possibles sont des caractères ASCII imprimables
(\underline{à l'exception de 0--3 et 99 centièmes de minute}).
\end{aprssideblock}

\chTenToc{Encodage de la vitesse et du cap}
\begin{aprssideblock}{Encodage de la\\vitesse et du cap}
La vitesse et le cap d'une station sont encodés sur 3 octets, désignés
\lit{SP+28}, \lit{DC+28} et \lit{SE+28}.

La vitesse est comprise entre 0 et 799 knots, et le cap entre 0 et 360 degrés
(0 degré représente un cap inconnu ou indéfini, et 360 degrés représente le
nord vrai).

La vitesse et le cap encodés sont répartis entre les trois octets comme suit :
\end{aprssideblock}

\begin{aprsbodyblock}
\begin{center}
\small\sffamily\renewcommand{\arraystretch}{1.50}
\setlength{\tabcolsep}{3pt}
\makebox[\linewidth][c]{\hspace*{-8pt}%
\begin{tabular}[t]{|>{\centering\arraybackslash}p{0.29\linewidth}|>{\centering\arraybackslash}p{0.34\linewidth}|>{\centering\arraybackslash}p{0.27\linewidth}|}\hline
\rowcolor{gray!20}\multicolumn{2}{|c|}{\chTenHead{Vitesse}}&\chTenHead{Cap}\\\hline
\chTenHead{Vitesse encodée\newline (centaines/dizaines de knots)}&
\chTenHead{Unités de vitesse encodées et\newline cap encodé (centaines de degrés)}&
\chTenHead{Cap encodé\newline (dizaines/unités)}\\\hline
\lit{SP+28}&\lit{DC+28}&\lit{SE+28}\\\hline
\end{tabular}%
\hspace*{8pt}}
\end{center}
\end{aprsbodyblock}

\chTenToc{Encodage SP+28}
\begin{aprssideblock}{Encodage SP+28}
L'octet \lit{SP+28} contient la vitesse encodée, en centaines/dizaines de
knots, conformément au tableau suivant :
\end{aprssideblock}

\begin{aprsbodyblock}
\sffamily\bfseries Encodage de la vitesse SP+28 (centaines/dizaines de knots)

\smallskip\footnotesize\sffamily\normalfont
\setlength{\tabcolsep}{2pt}\renewcommand{\arraystretch}{1.65}
\hspace*{-9.5pt}%
\begin{tabular}[t]{|>{\columncolor{gray!20}\bfseries\centering\arraybackslash}p{0.1495\linewidth}|>{\centering\arraybackslash}p{0.09\linewidth}|>{\columncolor{gray!20}\centering\arraybackslash}p{0.11\linewidth}|>{\centering\arraybackslash}p{0.065\linewidth}|>{\columncolor{gray!20}\centering\arraybackslash}p{0.065\linewidth}|}\hline
\rowcolor{gray!20}\rule{0pt}{16pt}\chTenHead{Vitesse\newline knots}&\multicolumn{2}{c|}{\chTenHead{ASCII Char}}&\multicolumn{2}{c|}{\chTenHead{SP+28}}\\\hline
0--9&\ttfamily l&\chTenBlack{0x1c}&108&28\\\hline
10--19&\ttfamily m&\chTenBlack{0x1d}&109&29\\\hline
20--29&\ttfamily n&\chTenBlack{0x1e}&110&30\\\hline
30--39&\ttfamily o&\chTenBlack{0x1f}&111&31\\\hline
40--49&\ttfamily p&&112&32\\\hline
50--59&\ttfamily q&\ttfamily !&113&33\\\hline
60--69&\ttfamily r&\aprsorigtexttt{"}&114&34\\\hline
70--79&\ttfamily s&\ttfamily \#&115&35\\\hline
80--89&\ttfamily t&\ttfamily \$&116&36\\\hline
90--99&\ttfamily u&\ttfamily \%&117&37\\\hline
100--109&\ttfamily v&\ttfamily \&&118&38\\\hline
110--119&\ttfamily w&\aprsorigtexttt{'}&119&39\\\hline
120--129&\ttfamily x&\ttfamily (&120&40\\\hline
130--139&\ttfamily y&\ttfamily )&121&41\\\hline
140--149&\ttfamily z&\ttfamily *&122&42\\\hline
150--159&\ttfamily \{&\ttfamily +&123&43\\\hline
160--169&\ttfamily |&\ttfamily ,&124&44\\\hline
170--179&\ttfamily \}&\ttfamily -&125&45\\\hline
180--189&\ttfamily \textasciitilde&\ttfamily .&126&46\\\hline
190--199&\chTenBlack{DEL}&\ttfamily /&127&47\\\hline
\end{tabular}\hspace{20.5pt}%
\begin{tabular}[t]{|>{\columncolor{gray!20}\bfseries\centering\arraybackslash}p{0.1495\linewidth}|>{\centering\arraybackslash}p{0.11\linewidth}|>{\columncolor{gray!20}\centering\arraybackslash}p{0.065\linewidth}|}\hline
\rowcolor{gray!20}\rule{0pt}{16pt}\chTenHead{Vitesse\newline knots}&\chTenHead{ASCII\newline Char}&\chTenHead{SP+28}\\\hline
200--209&\ttfamily 0&48\\\hline
210--219&\ttfamily 1&49\\\hline
220--229&\ttfamily 2&50\\\hline
230--239&\ttfamily 3&51\\\hline
240--249&\ttfamily 4&52\\\hline
250--259&\ttfamily 5&53\\\hline
260--269&\ttfamily 6&54\\\hline
270--279&\ttfamily 7&55\\\hline
280--289&\ttfamily 8&56\\\hline
290--299&\ttfamily 9&57\\\hline
300--310&\ttfamily :&58\\\hline
310--320&\ttfamily ;&59\\\hline
\ldots&&\\\hline
730--739&\ttfamily e&101\\\hline
740--749&\ttfamily f&102\\\hline
750--759&\ttfamily g&103\\\hline
760--769&\ttfamily h&104\\\hline
770--779&\ttfamily i&105\\\hline
780--789&\ttfamily j&106\\\hline
790--799&\ttfamily k&107\\\hline
\end{tabular}
\end{aprsbodyblock}

\begin{aprsbodyblock}
\textbf{Note :} les caractères ASCII représentés en blanc sur fond noir sont
des caractères non imprimables.

\textbf{Note :} pour les vitesses comprises entre 0 et 199 knots, deux schémas
d'encodage coexistent. Il y a donc deux colonnes pour le caractère ASCII et
deux colonnes pour les valeurs correspondantes de l'octet \lit{SP+28}.

Par exemple, pour une vitesse de 73 knots (c'est-à-dire comprise entre 70 et
79), l'octet \lit{SP+28} peut contenir soit \lit{s}, soit \lit{\#}, selon la
méthode d'encodage utilisée. Les deux sont également valides.

L'algorithme de décodage décrit plus loin traite indifféremment ces deux
schémas d'encodage.
\end{aprsbodyblock}

\chTenToc{Encodage DC+28}
\begin{aprssideblock}{Encodage DC+28}
L'octet \lit{DC+28} contient les unités de vitesse encodées, ainsi que le cap
encodé en centaines de degrés :
\end{aprssideblock}

\begin{center}
\begin{minipage}{0.95\textwidth}
\hspace*{14.035pt}\sffamily\bfseries Encodage vitesse / cap DC+28 (unités de knots/centaines de degrés)

\smallskip\footnotesize\sffamily\normalfont
\setlength{\tabcolsep}{2pt}\renewcommand{\arraystretch}{1.63}
\hspace*{14.535pt}%
\begin{tabular}[t]{|>{\columncolor{gray!20}\bfseries\centering\arraybackslash}p{38pt}|>{\columncolor{gray!20}\bfseries\centering\arraybackslash}p{45pt}|>{\centering\arraybackslash}p{23.5pt}|>{\columncolor{gray!20}\centering\arraybackslash}p{31.5pt}|>{\centering\arraybackslash}p{23.5pt}|>{\columncolor{gray!20}\centering\arraybackslash}p{32.5pt}|}\hline
\rowcolor{gray!20}\chTenHead{Knots\newline (units)}&\chTenHead{Course\newline (deg)}&\multicolumn{2}{c|}{\chTenHead{ASCII Char}}&\multicolumn{2}{c|}{\chTenHead{DC+28}}\\[-1.6pt]\hline
0&0--99&&\chTenBlack{0x1c}&32&28\\\hline
0&100--199&\ttfamily !&\chTenBlack{0x1d}&33&29\\\hline
0&200--299&\aprsorigtexttt{"}&\chTenBlack{0x1e}&34&30\\\hline
0&300--360&\ttfamily \#&\chTenBlack{0x1f}&35&31\\\hline
1&0--99&\ttfamily *&\ttfamily \&&42&38\\\hline
1&100--199&\ttfamily +&\aprsorigtexttt{'}&43&39\\\hline
1&200--299&\ttfamily ,&\ttfamily (&44&40\\\hline
1&300--360&\ttfamily -&\ttfamily )&45&41\\\hline
2&0--99&\ttfamily 4&\ttfamily 0&52&48\\\hline
2&100--199&\ttfamily 5&\ttfamily 1&53&49\\\hline
2&200--299&\ttfamily 6&\ttfamily 2&54&50\\\hline
2&300--360&\ttfamily 7&\ttfamily 3&55&51\\\hline
3&0--99&\ttfamily >&\ttfamily :&62&58\\\hline
3&100--199&\ttfamily ?&\ttfamily ;&63&59\\\hline
3&200--299&\ttfamily @&\ttfamily <&64&60\\\hline
3&300--360&\ttfamily A&\ttfamily =&65&61\\\hline
4&0--99&\ttfamily H&\ttfamily D&72&68\\\hline
4&100--199&\ttfamily I&\ttfamily E&73&69\\\hline
4&200--299&\ttfamily J&\ttfamily F&74&70\\\hline
4&300--360&\ttfamily K&\ttfamily G&75&71\\\hline
\end{tabular}\hspace{13.5pt}%
\begin{tabular}[t]{|>{\columncolor{gray!20}\bfseries\centering\arraybackslash}p{38pt}|>{\columncolor{gray!20}\bfseries\centering\arraybackslash}p{45pt}|>{\centering\arraybackslash}p{23.5pt}|>{\columncolor{gray!20}\centering\arraybackslash}p{23.5pt}|>{\centering\arraybackslash}p{23.5pt}|>{\columncolor{gray!20}\centering\arraybackslash}p{34.5pt}|}\hline
\rowcolor{gray!20}\chTenHead{Knots\newline (units)}&\chTenHead{Course\newline (deg)}&\multicolumn{2}{c|}{\chTenHead{ASCII Char}}&\multicolumn{2}{c|}{\chTenHead{DC+28}}\\[-1.6pt]\hline
5&0--99&\ttfamily R&\ttfamily N&82&78\\\hline
5&100--199&\ttfamily S&\ttfamily O&83&79\\\hline
5&200--299&\ttfamily T&\ttfamily P&84&80\\\hline
5&300--360&\ttfamily U&\ttfamily Q&85&81\\\hline
6&0--99&\ttfamily \textbackslash&\ttfamily X&92&88\\\hline
6&100--199&\ttfamily ]&\ttfamily Y&93&89\\\hline
6&200--299&\ttfamily \textasciicircum&\ttfamily Z&94&90\\\hline
6&300--360&\ttfamily \_&\ttfamily [&95&91\\\hline
7&0--99&\ttfamily f&\ttfamily b&102&98\\\hline
7&100--199&\ttfamily g&\ttfamily c&103&99\\\hline
7&200--299&\ttfamily h&\ttfamily d&104&100\\\hline
7&300--360&\ttfamily i&\ttfamily e&105&101\\\hline
8&0--99&\ttfamily p&\ttfamily l&112&108\\\hline
8&100--199&\ttfamily q&\ttfamily m&113&109\\\hline
8&200--299&\ttfamily r&\ttfamily n&114&110\\\hline
8&300--360&\ttfamily s&\ttfamily o&115&111\\\hline
9&0--99&\ttfamily z&\ttfamily v&122&118\\\hline
9&100--199&\ttfamily \{&\ttfamily w&123&119\\\hline
9&200--299&\ttfamily |&\ttfamily x&124&120\\\hline
9&300--360&\ttfamily \}&\ttfamily y&125&121\\\hline
\end{tabular}
\end{minipage}
\end{center}

\begin{aprsbodyblock}
\textbf{Note :} les caractères ASCII représentés en blanc sur fond noir sont
des caractères non imprimables.

\textbf{Note :} deux schémas d'encodage coexistent pour l'octet \lit{DC+28}.
Il y a donc deux colonnes pour le caractère ASCII et deux colonnes pour les
valeurs correspondantes de l'octet \lit{DC+28}.

Par exemple, pour une vitesse de 73 knots (c'est-à-dire unités = 3) et un cap
de 294 degrés (c'est-à-dire compris entre 200 et 299), l'octet \lit{DC+28}
peut contenir soit \lit{@}, soit \lit{<}, selon la méthode d'encodage utilisée.
Les deux sont également valides.

L'algorithme de décodage décrit plus loin traite indifféremment ces deux
schémas d'encodage.
\end{aprsbodyblock}

\chTenToc{Encodage SE+28}
\begin{aprssideblock}{Encodage SE+28}
L'octet \lit{SE+28} contient les dizaines et unités de degrés encodées du cap :

\par\smallskip
\sffamily\bfseries Encodage du cap SE+28 (dizaines/unités de degrés)

\smallskip\footnotesize\sffamily\normalfont
\setlength{\tabcolsep}{4.5pt}\renewcommand{\arraystretch}{1.68}
\hspace*{-3pt}%
\begin{tabular}[t]{|>{\columncolor{gray!20}\bfseries\centering\arraybackslash}p{33pt}|>{\centering\arraybackslash}p{26pt}|>{\centering\arraybackslash}p{26.5pt}|}\hline
\rowcolor{gray!20}\chTenHead{\shortstack{Course\\(deg)}}&\chTenHead{\shortstack{ASCII\\Char}}&\chTenHead{m+28}\\[7pt]\hline
0&\chTenBlack{0x1c}&28\\\hline
1&\chTenBlack{0x1d}&29\\\hline
2&\chTenBlack{0x1e}&30\\\hline
3&\chTenBlack{0x1f}&31\\\hline
4&&32\\\hline
5&!&33\\\hline
6&\aprsorigtexttt{"}&34\\\hline
7&\#&35\\\hline
8&\$&36\\\hline
9&\%&37\\\hline
10&\&&38\\\hline
11&\aprsorigtexttt{'}&39\\\hline
12&(&40\\\hline
13&)&41\\\hline
14&*&42\\\hline
\end{tabular}\hspace{28pt}%
\begin{tabular}[t]{|>{\columncolor{gray!20}\bfseries\centering\arraybackslash}p{33pt}|>{\centering\arraybackslash}p{26pt}|>{\centering\arraybackslash}p{26.5pt}|}\hline
\rowcolor{gray!20}\chTenHead{\shortstack{Long\\Mins}}&\chTenHead{\shortstack{ASCII\\Char}}&\chTenHead{m+28}\\[7pt]\hline
15&+&43\\\hline
16&,&44\\\hline
17&-&45\\\hline
18&.&46\\\hline
19&/&47\\\hline
\ldots&&\\\hline
91&w&119\\\hline
92&x&120\\\hline
93&y&121\\\hline
94&z&122\\\hline
95&\{&123\\\hline
96&|&124\\\hline
97&\}&125\\\hline
98&\textasciitilde&126\\\hline
99&\chTenBlack{DEL}&127\\\hline
\end{tabular}
\end{aprssideblock}

\chTenToc{Exemple d'encodage Mic-E de la vitesse et du cap}
\begin{aprssideblock}{Exemple d'encodage\\Mic-E de la vitesse\\et du cap}
Pour une vitesse de 86 knots et un cap de 194 degrés, l'encodage est le suivant :

\begin{description}
\setlength{\itemsep}{0.55em}\setlength{\leftmargin}{3.3em}
\item[\lit{SP+28} :] la vitesse est comprise entre 80 et 89 knots. D'après le
  tableau d'encodage SP+28, l'octet \lit{SP+28} peut être soit \lit{t}, soit
  \lit{\$}.
\item[\lit{DC+28} :] les unités de vitesse valent 6, et le cap est compris
  entre 100 et 199 degrés. D'après le tableau d'encodage DC+28, l'octet
  \lit{DC+28} peut être soit \lit{]}, soit \lit{Y}.
\item[\lit{SE+28} :] les dizaines et unités de degrés du cap valent 94.
  D'après le tableau d'encodage SE+28, l'octet \lit{SE+28} sera \lit{z}.
\end{description}
\end{aprssideblock}

\chTenToc{Décodage de la vitesse et du cap}
\begin{aprssideblock}{Décodage de la\\vitesse et du cap}
Pour décoder la vitesse et le cap :

\begin{description}
\setlength{\itemsep}{0.55em}\setlength{\leftmargin}{3.3em}
\item[\lit{SP+28} :] pour obtenir la vitesse en dizaines de knots, soustraire
  28 de la valeur \lit{SP+28} et multiplier par 10.
\item[\lit{DC+28} :] soustraire 28 de la valeur \lit{DC+28} et diviser le
  résultat par 10. Le quotient donne les unités de vitesse. Le reste donne le
  cap en centaines de degrés.
\item[\lit{SE+28} :] pour obtenir les dizaines et unités de degrés, soustraire
  28 de la valeur \lit{SE+28}.
\end{description}

Enfin, effectuer les ajustements suivants sur la vitesse et le cap :

\begin{itemize}\chTenItems
\item si la vitesse calculée est $\ge 800$ knots, soustraire 800 ;
\item si le cap calculé est $\ge 400$ degrés, soustraire 400.
\end{itemize}
\end{aprssideblock}

\chTenToc{Exemple de décodage des données du champ Information}
\begin{aprssideblock}{Exemple de décodage\\des données du champ\\Information}
Si les 9 premiers octets du champ Information contiennent
\lit{`(\_fn"Oj/}, et si la Destination Address indique que la station se trouve
dans l'hémisphère ouest avec un Longitude Offset de +100 degrés, les données
sont décodées comme suit :

\begin{itemize}\chTenItems
\item \lit{`} est l'APRS Data Type Identifier d'un paquet Mic-E contenant des
  données GNSS actuelles.
\item \lit{(} est l'octet \lit{d+28}. Le caractère \texttt{(} a la valeur 40
  en décimal. Soustraire 28 donne 12. Le Longitude Offset (dans la Destination
  Address) est de +100 degrés ; la longitude vaut donc $100 + 12 = 112$ degrés.
\item \lit{\_} est l'octet \lit{m+28}. Le caractère \texttt{\_} a la valeur
  95 en décimal. Soustraire 28 donne 67. Cette valeur étant $\ge 60$, soustraire
  60 donne une longitude de 7 minutes.
\item \lit{f} est l'octet \lit{h+28}. Le caractère \texttt{f} a la valeur 102
  en décimal. Soustraire 28 donne 74 centièmes de minute.
\end{itemize}

La longitude est donc de 112 degrés 7,74 minutes ouest.

La vitesse et le cap sont calculés comme suit :

\begin{itemize}\chTenItems
\item \lit{n} est l'octet \lit{SP+28}. Le caractère \texttt{n} a la valeur
  110 en décimal. Après soustraction de 28, le résultat vaut 82. Comme il est
  $\ge 80$, soustraire encore 80, ce qui laisse 2 dizaines de knots.
\item \lit{"} est l'octet \lit{DC+28}. Le caractère \texttt{"} a la valeur 34
  en décimal. Soustraire 28 donne 6. La division par 10 donne un quotient de 0
  (unités de vitesse). Ajouter au quotient (0) la première partie de la vitesse
  multipliée par 10 (soit 20) donne une vitesse finale calculée de 20 knots.

  Le reste de la division est 6. Soustraire 4 donne le cap en centaines de
  degrés, soit 2.
\item \lit{O} (lettre « O » majuscule) est l'octet \lit{SE+28}. Le caractère
  \texttt{O} a la valeur 79 en décimal. Soustraire 28 donne 51. Ajouter cette
  valeur au reste calculé ci-dessus, multiplié par 100 (soit 200), donne la
  valeur finale de 251 degrés pour le cap.
\end{itemize}

Les deux derniers caractères (\lit{j/}) représentent le symbole jeep de la
Primary Symbol Table.
\end{aprssideblock}

\chTenToc{Ambiguïté de position Mic-E}
\begin{aprssideblock}{Ambiguïté de\\position Mic-E}
Comme indiqué au chapitre 6 (Formats de temps et de position), une station peut
réduire la précision de sa position en introduisant une ambiguïté de position.
Cela est également possible dans le format de données Mic-E.

L'ambiguïté de position est spécifiée pour la latitude (dans la Destination
Address). Le même degré d'ambiguïté s'applique alors également à la longitude.

Par exemple, si la Destination Address est \lit{T4SQZZ}, les deux derniers
chiffres de la latitude sont ambigus (représentés par \lit{ZZ}). Alors, si les
données de longitude du champ Information sont \lit{(\_f}, comme dans l'exemple
ci-dessus, les deux derniers chiffres de la longitude calculée sont ignorés :
la longitude vaut donc 112 degrés 7 minutes.
\end{aprssideblock}

\chTenToc{Données de télémétrie Mic-E (obsolètes)}
\begin{aprssideblock}{Données de télémétrie\\Mic-E (obsolètes)}
Ce format est désormais obsolète et remplacé par le format de télémétrie
Base 91 Comment décrit dans le chapitre Télémétrie.

Si l'octet qui suit le Symbol Table Identifier est l'un des caractères
Telemetry Flag (\lit{`}, \lit{'}, ou \texttt{0x1d}), des données de
télémétrie suivent :
\end{aprssideblock}

\begin{center}
\small\sffamily\setlength{\tabcolsep}{4.5pt}\renewcommand{\arraystretch}{1.45}
\makebox[\textwidth][c]{\hspace*{5pt}%
\begin{tabular}[t]{|>{\centering\arraybackslash}p{0.09\textwidth}|*{5}{>{\centering\arraybackslash}p{0.055\textwidth}|}}\hline
\multicolumn{6}{|l|}{\sffamily\bfseries\itshape Données de télémétrie Mic-E optionnelles (obsolètes)}\\\hline
\rowcolor{gray!20}\chTenHead{Telemetry\newline Flag}&\multicolumn{5}{c|}{\chTenHead{Canaux de données de télémétrie}}\\\hline
F&Ch 1&Ch 2&Ch 3&Ch 4&Ch 5\\\hline
\makebox[0pt][r]{Bytes:\hspace{1.2em}}\makebox[\linewidth][c]{1}&1/2&1/2&1/2&1/2&1/2\\\hline
\end{tabular}%
\hspace*{-5pt}}
\end{center}

\begin{aprsbodyblock}
Le Telemetry Flag F est l'un des suivants :

\begin{tabular}[t]{@{}>{\ttfamily}p{0.07\linewidth}p{0.86\linewidth}@{}}
\lit{`} & 2 valeurs de télémétrie hexadécimales imprimables suivent (canaux 1
et 3) ; \\
\lit{'} & 5 valeurs de télémétrie hexadécimales imprimables suivent ; \\
\texttt{0x1d} & 5 valeurs de télémétrie binaires suivent (unités beta Rev. 0
uniquement).
\end{tabular}

Si F vaut \lit{`} ou \lit{'}, chaque canal nécessite 2 octets, contenant une
représentation hexadécimale imprimable sur 2 chiffres d'une valeur comprise
entre 0 et 255. Par exemple, 254 est représenté par \texttt{FE}.

Si F vaut \texttt{0x1d}, chaque canal nécessite un octet contenant une
valeur binaire de 8 bits.

Par exemple, si les données de télémétrie sont \lit{'7200007100}, le caractère
\lit{'} indique que 5 octets de télémétrie suivent, codés en hexadécimal :

\begin{tabular}[t]{@{}l@{}}
\texttt{0x72 = 114 decimal}\\
\texttt{0x00 = 0 decimal}\\
\texttt{0x00 = 0 decimal}\\
\texttt{0x71 = 113 decimal}\\
\texttt{0x00 = 0 decimal}
\end{tabular}
\end{aprsbodyblock}

\chTenToc{Texte de statut Mic-E}
\begin{aprssideblock}{Texte de statut Mic-E}
Le paquet peut inclure du texte de statut Mic-E. Le texte de statut peut avoir
toute longueur tenant dans le reste du champ Information. (Pourquoi ne pas
simplement l'appeler « commentaire », par cohérence avec les autres formats de
rapports de position et d'objets ?)

Le texte de statut Mic-E ne doit pas commencer par \lit{`}, \lit{'}, ou
\texttt{0x1d}, faute de quoi il serait confondu avec des données de
télémétrie [désormais obsolètes].

Il est possible d'inclure une position au format APRS standard dans le champ de
texte de statut Mic-E. \emph{\textcolor{red}{Hein ???}} Une position appropriée amènera le
logiciel d'affichage APRS à remplacer toutes les données de position encodées
par le Mic-E. Cela est utile lorsqu'un Mic-E est utilisé sans récepteur GNSS.
\emph{\textcolor{red}{Hein ???}}

Le champ de texte Mic-E peut aussi contenir n'importe quel champ commentaire de
position normal, tel que PHG.

Pour se distinguer du dispositif MIC-E d'origine, les premières radios Kenwood
inséraient automatiquement un code de préfixe au début de la chaîne de texte de
statut (c'est-à-dire au 10e caractère du champ Information) :

\begin{tabular}[t]{@{}ll@{}}
Kenwood TH-D7: & \lit{>}\\
Kenwood TM-D700: & \lit{]}
\end{tabular}

Les modèles Kenwood ultérieurs ont conservé le même préfixe, pour indiquer un
portatif ou un mobile, puis ont ajouté un caractère de suffixe unique au texte :

\begin{tabular}[t]{@{}ll@{}}
Kenwood TH-D72: & \lit{>} \ldots \lit{=}\\
Kenwood TH-D74: & \lit{>} \ldots \lit{\textasciicircum}\\
Kenwood TH-D75: & \lit{>} \ldots \lit{\&}\\
Kenwood TM-D710: & \lit{]} \ldots \lit{=}
\end{tabular}

Les dispositifs d'autres fabricants utilisaient un caractère de préfixe
différent :
\lit{`} indique qu'un système est capable de messagerie, ou
\lit{'} indique qu'un système n'est pas capable de messagerie (par exemple un
tracker). Un suffixe de deux caractères indique le fabricant et le modèle.
Exemples :

\begin{tabular}[t]{@{}ll@{}}
Yaesu FT2D: & \lit{`} \ldots \lit{\_(} (messagerie)\\
Yaesu FT3D: & \lit{`} \ldots \lit{\_0}\\
Yaesu FT5D: & \lit{`} \ldots \lit{\_3}\\
Byonics TinyTrack 3: & \lit{'} \ldots \lit{|3} (sans messagerie)\\
Byonics TinyTrack 4: & \lit{'} \ldots \lit{|4}
\end{tabular}

Les applications d'affichage APRS doivent supprimer tout préfixe et suffixe
supplémentaire avant d'afficher le texte. Si vous voyez quelque chose comme
« \lit{\_3} » à la fin du texte de statut (commentaire), c'est parce que
l'application n'a pas supprimé le type de dispositif comme elle aurait dû.

Plutôt que d'utiliser des identifiants de dispositif codés en dur, il est
fortement recommandé que les applications lisent un fichier à l'exécution.
Lorsque de nouveaux modèles sont ajoutés, seul ce fichier doit être mis à jour,
plutôt que de modifier le code et de publier une nouvelle version. Des fichiers
lisibles par machine, avec les derniers identifiants de dispositif, sont
disponibles à l'adresse \url{https://github.com/aprsorg/aprs-deviceid}.
\end{aprssideblock}

\chTenToc{Locator Maidenhead dans le champ de texte de statut Mic-E}
\begin{aprssideblock}{Locator Maidenhead\\dans le champ de texte\\de statut Mic-E}
Le champ de texte de statut Mic-E peut contenir un locator Maidenhead.

Si le locator est suivi d'un commentaire en texte clair, le premier caractère
du texte doit être une espace. Par exemple :

\begin{tabular}[t]{@{}ll@{}}
\aprsorigtexttt{IO91SX/G}\chTenLitSpace\aprsorigtexttt{Hello}\chTenSpace\aprsorigtexttt{world} & (depuis un Mic-E ou PIC-E)\\
\aprsorigtexttt{>IO91SX/G}\chTenLitSpace\aprsorigtexttt{Hello}\chTenSpace\aprsorigtexttt{world} & (depuis un Kenwood TH-D7)\\
\aprsorigtexttt{]IO91SX/G}\chTenLitSpace\aprsorigtexttt{Hello}\chTenSpace\aprsorigtexttt{world} & (depuis un Kenwood TM-D700)
\end{tabular}

(\lit{/G} est le symbole de grid locator).
\end{aprssideblock}

\chTenToc{Altitude dans le champ de texte de statut Mic-E}
\begin{aprssideblock}{Altitude dans le champ\\de texte de statut\\Mic-E}
Le champ de texte de statut Mic-E peut contenir l'altitude de la station.
L'altitude est exprimée sous la forme \aprsorigtexttt{xxx}\lit{\}}, où
\aprsorigtexttt{xxx} est en mètres
par rapport à 10 km sous le niveau moyen de la mer (l'océan le plus profond),
en base 91.

Par exemple, pour calculer les caractères \aprsorigtexttt{xxx} pour une altitude de
200 pieds :

\begin{tabular}[t]{@{}l@{}}
200 pieds = 61 mètres = 10061 mètres par rapport au datum\\
\(10061 / 91^2 = {}\)\lit{1}, reste 1780\\
\(1780 / 91\phantom{^2} = {}\)\lit{19}, reste \lit{51}
\end{tabular}

L'ajout de 33 à chacune des valeurs surlignées donne 34, 52 et 84 pour les
codes ASCII de \aprsorigtexttt{xxx}.

La chaîne d'altitude de 4 caractères est donc \lit{"4T\}}.

Quelques exemples :

\begin{tabular}[t]{@{}ll@{}}
\aprsorigtexttt{"4T\}}\\
\aprsorigtexttt{>"4T\}}\\
\aprsorigtexttt{]"4T\}}
\end{tabular}

L'altitude optionnelle doit être placée en premier après l'octet de type
Mic-E.

Noter que cela vient après tout caractère de préfixe d'identifiant de
dispositif.
\end{aprssideblock}

\chTenToc{Données Mic-E dans les réseaux non APRS}
\begin{aprssideblock}{Données Mic-E dans\\les réseaux non APRS}
Certaines parties du champ Information AX.25 Mic-E peuvent contenir des
données binaires (c'est-à-dire des caractères ASCII non imprimables). Si un
tel paquet reste limité au réseau APRS, cela ne devrait poser aucune
difficulté.

Toutefois, si le paquet doit être transmis par un réseau qui ne préserve pas
fiablement les données binaires (Internet, par exemple), il est nécessaire de
convertir les données dans un format qui les préservera.

En outre, si le paquet revient ensuite sur le réseau APRS, il est alors
nécessaire de reconvertir les données dans leur format d'origine.
\end{aprssideblock}

\let\chTenToc\relax
\let\chTenHead\relax
\let\chTenGray\relax
\let\chTenBlack\relax
\let\chTenItems\relax
\endgroup
